\scanpage{14}
il est immédiat \unsure{à priori|a priori} que \unsure{tout} ss-schéma
$\deleted{\unsure{A}}$ en groupes \inserted{fermé} $B'$ de $A'=A_K$ est
la \unsure{contraction|restriction} d'un \deleted{\unsure{ab}} ss-schéma abélien
$B$ de $A$. Il s'ensuit, si \unsure{maintt|maintenant}
$A$ est un schéma abélien \deleted{\unsure{pa|projectif}} \deleted{\unsure{quel}} sur un
schéma loc. noeth $S$ \uline{de car. $0$}, \deleted{\unsure{que|tel}} tel
\deleted{le schéma} que $\uuline{\mathrm{Hilb}}_{\inserted{A/S}}\deleted{(A/S)}$ existe \inserted{et est loc. de type fini}
(p. ex. \deleted{\unsure{?}} \deleted{\unsure{que}} \inserted{si} $\deleted{\unsure{B}}A$ projectif sur $S$) alors
le schéma $\unsure{\mathscr{M}|\mathcal{M}}$ sur $S$ de page $2$
\deleted{est tel que} qui classifie les \uline{sous-}schémas \uline{abéliens}
\deleted{(resp. \unsure{séparables})} est
\deleted{tel que ses parties \inserted{\unsure{fermées}} de type fini}
\uline{loc. propre sur $S$}. D'ailleurs, en
toute caractéristique, le $S$-préschéma $N_1$
qui classifie les ss-schémas \deleted{\uline{séparables}} \inserted{en groupes \uline{fermés}} de
$A$ est \deleted{\unsure{localement}} \unsure{sur|\deleted{sur}} \inserted{\uline{loc. non ramifié}} \inserted{\unsure{!|1}} $S$. Donc
En caractéristique $0$ en combinant les,
deux \deleted{\unsure{énoncés}} propriétés, on voit ceci:
% archive sheet number: 13
